% SPDX-License-Identifier: MIT
% Copyright (c) 2026 Evgeny Baulin
% Optimization Methods, Seminar 07: one-page cheat sheet (English).
% Compile with: latexmk -xelatex -interaction=nonstopmode -halt-on-error Cheatsheet_en.tex
\RequirePackage{fix-cm}
\documentclass[10pt,a4paper]{article}

\usepackage{amsmath,amssymb}
\usepackage{fontspec}
\setmainfont{cmunrm.otf}[BoldFont=cmunbx.otf, ItalicFont=cmunti.otf, BoldItalicFont=cmunbi.otf]
\setsansfont{cmunss.otf}[BoldFont=cmunsx.otf, ItalicFont=cmunsi.otf, BoldItalicFont=cmunso.otf]
\setmonofont{cmuntt.otf}[BoldFont=cmuntb.otf, ItalicFont=cmunit.otf, BoldItalicFont=cmuntx.otf]
\usepackage{polyglossia}
\setdefaultlanguage{english}

\usepackage[a4paper,left=8mm,right=8mm,top=7mm,bottom=11mm,footskip=5.5mm]{geometry}
\usepackage{xcolor}
\usepackage{multicol}
\usepackage{enumitem}
\usepackage{ragged2e}
\usepackage{fancyhdr}
\usepackage[colorlinks=true,linkcolor=hseblue,urlcolor=hseblue]{hyperref}

\definecolor{hseblue}{RGB}{16,45,105}
\definecolor{hseaccent}{RGB}{198,40,40}
\definecolor{hsegray}{RGB}{108,114,126}
\definecolor{hselight}{RGB}{240,243,248}

\hypersetup{pdftitle={Seminar 07 cheat sheet: Pre-colloquium Practice}, pdfauthor={Evgeny Baulin}}

\pagestyle{fancy}
\fancyhf{}
\renewcommand{\headrulewidth}{0pt}
\fancyfoot[L]{\scriptsize\color{hsegray}\textcopyright\ 2026 Evgeny Baulin, \href{https://t.me/tarakan_tuc}{t.me/tarakan\_tuc}}
\fancyfoot[R]{\scriptsize\color{hsegray}MIT License}

\setlength{\parindent}{0pt}
\setlength{\parskip}{0pt}
\setlength{\columnsep}{4.5mm}
\setlength{\columnseprule}{0.3pt}
\renewcommand{\columnseprulecolor}{\color{hselight}}
\setlength{\emergencystretch}{2em}
\setlength{\RaggedRightRightskip}{0pt plus 4em}
\binoppenalty=2000
\relpenalty=500
\setlength{\abovedisplayskip}{2pt}
\setlength{\belowdisplayskip}{2pt}
\setlist[itemize]{leftmargin=7pt, labelsep=3pt, itemsep=1.2pt, topsep=1pt, parsep=0pt, partopsep=0pt,
  label={\color{hseaccent}\scriptsize\textbullet}}

\makeatletter
\newcommand{\cssec}[1]{\par\addvspace{3.5pt}\noindent
  \colorbox{hselight}{\parbox{\dimexpr\linewidth-2\fboxsep}{\RaggedRight\hyphenpenalty=10000\color{hseblue}\bfseries\strut #1}}%
  \par\nobreak\vspace{1pt}\@afterheading}
\makeatother
\clubpenalty=10000
\widowpenalty=10000

\begin{document}
\thispagestyle{fancy}

\noindent
\begin{minipage}[b]{0.6\textwidth}
  \RaggedRight\hyphenpenalty=10000
  {\color{hseaccent}\bfseries\small Seminar 07 cheat sheet}\\[1pt]
  {\color{hseblue}\bfseries\Large Pre-colloquium\\ Practice}
\end{minipage}\hfill
\begin{minipage}[b]{0.39\textwidth}
  \raggedleft\footnotesize
  {\bfseries Optimization Methods}\\
  HSE University, Faculty of Computer Science, 2026--2027\\
  Course lecturer: Andrey Ignatov\\
  Author: Evgeny Baulin, Telegram \href{https://t.me/tarakan_tuc}{@tarakan\_tuc}
\end{minipage}

\vspace{2pt}
{\color{hseblue}\rule{\textwidth}{0.8pt}}
\vspace{-4pt}

\fontsize{9}{11.4}\selectfont
\begin{multicols}{3}
\RaggedRight
\cssec{Definitions you must know}
\begin{itemize}
  \item \textbf{Convex set} $C$: $tx+(1-t)y\in C$ for every $x,y\in C$ and $t\in[0,1]$.
  \item \textbf{Convex function} on convex $C$: $f(tx+(1-t)y)\le tf(x)+(1-t)f(y)$.
  \item \textbf{Strongly convex} with $\mu>0$: $f(y)\ge f(x)+\nabla f(x)^\top(y-x)+\tfrac{\mu}{2}\|y-x\|^2$.
  \item \textbf{$L$-smooth}: $\|\nabla f(y)-\nabla f(x)\|\le L\|y-x\|$ for every $x,y$.
  \item \textbf{Gradient}: $\nabla f(x)=(\partial_1f,\dots,\partial_nf)^\top$. \textbf{Hessian}: $(\nabla^2 f)_{ij}=\partial_i\partial_j f$, symmetric.
  \item \textbf{Stationary point}: $\nabla f(x^*)=0$. \textbf{Minimum, maximum, saddle} by signs of eigenvalues of $\nabla^2 f(x^*)$.
\end{itemize}

\cssec{Problem 3.1: a quadratic form}
\begin{itemize}
  \item $f(x,y)=ax^2+bxy+cy^2+dx+ey$: $\nabla f=(2ax+by+d,\ bx+2cy+e)$.
  \item $H=\begin{pmatrix}2a&b\\b&2c\end{pmatrix}$. Solve $Hx^*=-(d,e)^\top$.
  \item Eigenvalues: $\lambda^2-\operatorname{tr}(H)\lambda+\det(H)=0$.
  \item Both $>0$: minimum; both $<0$: maximum; different signs: saddle.
  \item $\mu=\lambda_{\min}(H)$, $L=\lambda_{\max}(H)$ when $H\succ0$.
\end{itemize}

\cssec{Problem 3.2: gradient descent}
\begin{itemize}
  \item $f=\tfrac{1}{2}x^\top Qx$, $Q\succ0$: $x_{k+1}=(I-\alpha Q)x_k$.
  \item Safe: $0<\alpha<2/L$. Border: $\alpha=2/L$ stalls; $\alpha>2/L$ diverges.
  \item In the eigenbasis: $e_{k,i}=(1-\alpha\lambda_i)^k e_{0,i}$.
  \item $\alpha=1/L$: steepest axis in one step; slow axis by $1-\mu/L$.
  \item Best fixed step $\alpha^*=2/(\mu+L)$: rate $(\kappa-1)/(\kappa+1)$ per step.
\end{itemize}

\cssec{Two short proofs}
\begin{itemize}
  \item \textbf{Convex combination.} $f,g$ convex, $\alpha,\beta\ge0$: $h=\alpha f+\beta g$ is convex by linearity in $f,g$.
  \item \textbf{Strict minimum.} $\nabla f(x^*)=0$, $\nabla^2 f(x^*)\succ0$: Taylor gives $f(x^*+h)=f(x^*)+\tfrac{1}{2}h^\top\nabla^2 f(x^*)h+o(\|h\|^2)>f(x^*)$ for small $h\ne0$.
\end{itemize}

\cssec{What to watch for}
\begin{itemize}
  \item $H$ has $2a$ and $2c$ on the diagonal, not $a$ and $c$.
  \item Zero eigenvalue: the second-order test fails, higher-order terms needed.
  \item $\alpha=1/L$ is not the best step; $\alpha=2/(\mu+L)$ is.
\end{itemize}

\end{multicols}
\end{document}

